\documentclass[a4paper,11pt]{article}
\usepackage[utf8]{inputenc}
\usepackage[T1]{fontenc}
\usepackage[french]{babel}
\usepackage{lmodern}
\usepackage{amsmath,amssymb,amsthm}
\usepackage{mathrsfs}
\usepackage{enumitem}
\usepackage{multicol}

\setlist[enumerate]{label=\textbf{\textcolor{Couleur8}{\arabic*.}}, left=1em}

% Si vous voulez aussi modifier les listes enumerate imbriquées :
\setlist[enumerate,2]{label=\textcolor{Couleur8}{\alph*)}}
\setlist[enumerate,3]{label=\textcolor{Couleur8}{\roman*)}}
\setlist[enumerate,4]{label=\textcolor{Couleur8}{\Alph*)}}
\usepackage{graphicx}
\usepackage{xcolor}
\usepackage{tcolorbox}
\usepackage{geometry}
\usepackage{fancyhdr}
\usepackage{titlesec}
\usepackage{cancel}
\usepackage{ifthen}
\usepackage{tikz}
\usetikzlibrary{arrows.meta}
\usepackage{tkz-tab}
\usepackage{eurosym}

% OPTION PROF/ÉLÈVE
% Décommentez UNE des deux lignes suivantes :
\newboolean{prof}\setboolean{prof}{true}   % Version professeur (avec corrections des devoirs)
%\newboolean{prof}\setboolean{prof}{false}  % Version élève (sans corrections des devoirs)

\definecolor{Couleur1}{HTML}{4A5D7A} %Th
\definecolor{Couleur2}{HTML}{A5B5D6} %Prop
\definecolor{Couleur3}{HTML}{A8C68A} %Méthode
\definecolor{Couleur6}{HTML}{8A9CC1} %Def
\definecolor{Couleur7}{HTML}{E6B459} %Rq Voc
\definecolor{Couleur8}{HTML}{D36740} %Titres
\definecolor{correction}{HTML}{D36740} % Couleur pour les corrections
\definecolor{coopmathscorr}{HTML}{F15929} % Couleur corrections coopmaths

% Configuration des marges
\geometry{
	top=2cm,
	bottom=2cm,
	inner=1.5cm,
	outer=1.5cm,
}

% Police
\renewcommand{\familydefault}{\sfdefault}

% Compteurs
\newcounter{seancenum}
\newcounter{seancenumD}
\setcounter{seancenum}{1}
\setcounter{seancenumD}{1}

% Configuration des en-têtes et pieds de page
\pagestyle{fancy}
\fancyhf{}
\ifthenelse{\boolean{prof}}{
    \fancyhead[L]{\textcolor{Couleur8}{\textbf{Corrections Automatismes - Classe de seconde - Version Professeur}}}
}{
    \fancyhead[L]{\textcolor{Couleur8}{\textbf{Corrections Automatismes - Séance 29 - Classe de seconde}}}
}
\fancyhead[R]{\textcolor{Couleur8}{\textbf{2025-2026}}}
\fancyfoot[L]{\textcolor{Couleur8}{Mme Bertail et Mme Pinto}}
\fancyfoot[R]{\textcolor{Couleur8}{\thepage}}
\renewcommand{\headrulewidth}{1pt}
\renewcommand{\headrule}{\hbox to\headwidth{\color{Couleur8}\leaders\hrule height \headrulewidth\hfill}}

% Environnements personnalisés
\newtcolorbox{seance}[1]{
    colback=Couleur6!10,
    colframe=Couleur1!65,
    fonttitle=\bfseries\large,
    title=Correction Séance \theseancenum #1,
    left=5mm,
    right=5mm,
    top=3mm,
    bottom=3mm,
    arc=0mm,
    after=\stepcounter{seancenum}
}

\newtcolorbox{seanceD}[1]{
    colback=Couleur6!5,
    colframe=Couleur6!45,
    fonttitle=\bfseries\large,
    title=Correction Devoirs - Séance \theseancenumD #1,
    left=5mm,
    right=5mm,
    top=3mm,
    bottom=3mm,
    arc=0mm,
    after=\stepcounter{seancenumD}
}

\begin{document}

\setcounter{seancenum}{29}
\newpage
\begin{seance}{} %29

\begin{enumerate}
\item L'ordonnée de ce point est $0$ puisque le point d'intersection se situe sur l'axe des abscisses.\\
      Son abscisse est donc donnée par la solution de l'équation  $3x+6=0$, c'est-à-dire $x=-2$.
    \\Les coordonnées de ce   point sont donc : $({\color[HTML]{F15929}\boldsymbol{-2\,;\,0}})$.\\La bonne réponse est la réponse {\bfseries \color[HTML]{F15929}C}.

\item On commence par déterminer l'équation réduite de la droite $(D)$.\\
    Son coefficient directeur est $-\dfrac{1}{5}$ et son ordonnée à l'origine est $1$.\\
    L'équation réduite de $(D)$ est : $y=-\dfrac{1}{5}x+1$.\\
    L'ordonnée du point de $(D)$ d'abscisse $60$ est donnée par : $y=-\dfrac{1}{5}\times 60+1$, ce qui donne\\
    $y=\dfrac{-60}{5}+1=-12+1=-11$.\\
    L'ordonnée du point de $(D)$ d'abscisse $60$ est $-11$.\\On a $-14<-11$, donc {\bfseries \color[HTML]{F15929}le point $A$ est situé en dessous de la droite $(D)$}.\\La bonne réponse est la réponse {\bfseries \color[HTML]{F15929}C}.

\item Pour vérifier si un point $(x\,;\,y)$ appartient à la courbe d'équation $y=\dfrac{3}{x}$, on vérifie si les coordonnées satisfont l'équation.\\$\bullet$ Pour le point $Q$ :\\On calcule $\dfrac{3}{\dfrac{6}{5}}=3\times \dfrac{5}{6}=\dfrac{5}{2}$.\\On a $y_{Q}=\dfrac{18}{5}\neq \dfrac{5}{2}$ donc {\bfseries \color[HTML]{F15929}$Q$ n'appartient pas à la courbe}.\\$\bullet$ Pour le point $W$ :\\On calcule $\dfrac{3}{\dfrac{4}{5}}=3\times \dfrac{5}{4}=\dfrac{15}{4}$.\\On a $y_{W}=\dfrac{15}{4}$ donc $W$ appartient à la courbe.\\$\bullet$ Pour le point $X$ :\\On calcule $\dfrac{3}{\dfrac{2}{3}}=3\times \dfrac{3}{2}=\dfrac{9}{2}$.\\On a $y_{X}=\dfrac{9}{2}$ donc $X$ appartient à la courbe.\\$\bullet$ Pour le point $V$ :\\On calcule $\dfrac{3}{\dfrac{12}{7}}=3\times \dfrac{7}{12}=\dfrac{7}{4}$.\\On a $y_{V}=\dfrac{7}{4}$ donc $V$ appartient à la courbe.

\medskip
La bonne réponse est la réponse {\bfseries \color[HTML]{F15929}B}.

\item L'équation est de la forme $X^2=k$ avec $X=(x+3)$ et $k=10$.\\
        Comme $k>0$, les solutions de $(x+3)^2=10$ sont données par les solutions de chacune des équations :         $x+3=-\sqrt{10}$ et $x+3=\sqrt{10}$.\\
        $x+3=-\sqrt{10}$ a pour solution $-\sqrt{10}-3$ et 
        $x+3=\sqrt{10}$ a pour solution $\sqrt{10}-3$.\\
        Ainsi, l'ensemble des solutions de l'équation est $S={\color[HTML]{F15929}\boldsymbol{\{-\sqrt{10}-3\,;\,\sqrt{10}-3\}}}$.
         \\La bonne réponse est la réponse {\bfseries \color[HTML]{F15929}C}.

\item Dans une $1$ heure, il y a $5\times 12$ minutes.\\
      Ainsi, en $1$ heure, elle parcourt $12$ fois plus de distance, soit $0{,}5\times 12=6\text{ km}$.\\
      Sa vitesse moyenne est donc ${\color[HTML]{F15929}\boldsymbol{6}}\text{ km/h}$. \\La bonne réponse est la réponse {\bfseries \color[HTML]{F15929}A}.

\item La distance Paris-Lyon est d'environ $465$ ${\color[HTML]{F15929}\boldsymbol{\textbf{km}}}$.\\La bonne réponse est la réponse {\bfseries \color[HTML]{F15929}C}.

\end{enumerate}

\end{seance}

\end{document}